\begin{lemma}
Let $G$ be a finite simple graph of maximum degree at most $\Delta$.  Then
there is a finite simple $\Delta$-regular graph $G'$ and an injective map
$\phi:V(G)\to V(G')$ such that, for all old vertices $u,v\in V(G)$,
\[
uv\in E(G) \quad\Longleftrightarrow\quad
\phi(u)\phi(v)\in E(G').
\]
Thus the image of $G$ is an induced subgraph of $G'$.
\end{lemma}
\begin{proof}
If $\Delta=0$, the maximum-degree hypothesis says that $G$ has no edges.
Take $G'=G$ and $\phi$ to be the identity; every vertex then has degree
$0$, and the induced-subgraph assertion is immediate.

Assume $\Delta>0$.  Start with a disjoint copy of $G$ and call its vertices
old.  For each old vertex $v$ of degree $d_G(v)$, create
\[
\Delta-d_G(v)
\]
formal stubs incident with $v$.  Filling each stub by one new edge will make
the old vertex have degree exactly $\Delta$.  Let $s$ be the total number of
old stubs.  If $s$ is odd, add one fresh isolated vertex and give it
$\Delta$ stubs.  Since an odd $s$ can occur only when $\Delta$ is odd by the
handshaking parity congruence
\[
s=\Delta |V(G)|-2|E(G)|\equiv \Delta |V(G)|\pmod 2,
\]
the number $s+\Delta$ is even.  If $s$ is already even, add no parity
vertex.  We now have an even finite set of stubs.

Pair the stubs arbitrarily.  For each pair of stubs, introduce a fresh copy
of the complete graph $K_{\Delta+1}$ and choose one edge $ab$ inside that
copy.  Delete the edge $ab$.  If the two paired stubs belong to vertices
$x$ and $y$, add the edges $xa$ and $yb$.  The two vertices $a$ and $b$ had
degree $\Delta-1$ after deleting $ab$, and each regains one incident edge;
all other vertices of the complete-graph block keep degree $\Delta$.
Meanwhile each stub-owner receives exactly one new incident edge.  If
$x=y$, the two new neighbors $a$ and $b$ are distinct, so no loop is formed.
If $x\ne y$, both new endpoints are fresh and lie in a block not used for
any other stub-pair, so no duplicate edge is formed.  No edge is ever added
between two old vertices.

After all stub pairs are filled, every old vertex has gained exactly the
number of missing incident edges, the possible parity vertex has also had
all of its stubs filled, and every vertex in every fresh
$K_{\Delta+1}$ block has degree $\Delta$.  The resulting graph is finite
because it is obtained from finitely many stubs and finitely many finite
blocks.  It is simple by the no-loop and no-duplicate checks above.  Since
no old-old edge was added or removed, two old vertices are adjacent in the
final graph exactly when they were adjacent in $G$.  Hence the inclusion of
old vertices is injective and identifies $G$ with an induced subgraph of the
finite $\Delta$-regular graph just constructed.
\end{proof}
